\addtocontents{toc}{\protect\Needspace{3\baselineskip}}
\originalchapter{黑洞热力学与全息动机}\label{ch:holography}
\originalsection{欧氏正则性与温度}
黑洞因果结构与 Schwarzschild 几何承接广义相对论卷. 本章需要的是其热性质. 非极端静态视界附近 $\dd s^2=-f(r)\dd t^2+\dd r^2/h(r)+\cdots$, $f,h$ 在 $r_0$ 有一阶零点. 令 $t=-\ii t_E$、$r-r_0=h'(r_0)\rho^2/4$, 则径向时间部分近似
$\dd s_E^2=\dd\rho^2+[f'(r_0)h'(r_0)/4]\rho^2\dd t_E^2$.
它是极坐标平面, 无锥尖需角周期 $2\pi$, 所以
\begin{equation}\label{eq:Hawking}
T_H=\frac{\sqrt{f'(r_0)h'(r_0)}}{4\pi}.
\end{equation}
时间坐标需按远处物理时钟归一化; 任意重缩放 $t$ 会改变式中温度. 极端双零点不适用这一步极坐标推导.

Schwarzschild 黑洞 $r_s=2GM$ 给 $T_H=1/(8\pi GM)$. 一阶定律 $\dd M=T_H\dd S$ 积分得 $S=A/(4G)$, 选零面积零熵的常数. 恢复单位为 $S=\kb c^3A/(4G\hbar)$. 熵正比面积而不是体积, 提示量子引力的自由度计数与普通局域场的紫外模计数不同. 这一积分使用 Hawking 温度和热力学一阶定律, 不构成任意黑洞微观态计数的证明.

\originalsection{熵界与面积}
半径 $R$ 的普通辐射气体在重力坍缩前的最大能量数量级 $E\lesssim R/(2G)$. 自由辐射 $S\propto Z^{1/4}V^{1/4}E^{3/4}$, 代入得 $S_{\rm gas}\propto Z^{1/4}(R/l_P)^{3/2}$, 比 $S_{\rm BH}\propto(R/l_P)^2$ 少. 对宏观大小和通常粒子种数, 普通热气体远未用尽黑洞熵. 但大量种数和极高密度会改变有效理论截止, 不能用这个估计证明无条件、所有引力体系都满足一个简单静态熵界.

全息原则是以边界面积刻画可用自由度的思想. AdS/CFT 提供具体的边界量子系统候选, 并把边界尺度与体径向方向联系起来. Weinberg--Witten 定理的条件包括 Lorentz 协变守恒电流或应力张量; 规范场与引力可通过不满足这些条件避开简单无禁结论. 本卷不展开该定理完整证明及所有逃逸机制, 它作为研究动机列在待补范围.
